\documentclass[a4paper,12pt]{article}

\usepackage{titling}
\usepackage{graphicx}

% Define custom title format
\renewcommand{\maketitle}{
  \begin{center}
    {\LARGE\textbf{Comparison of Machine Learning Algorithms}}\par
    {\Large Project Plan}\par
    \vspace{10pt}
    Nathan Kurien\par
    CS3821: BSc Final Project\par
    \hrule % Horizontal line
    \vspace{20pt}
    \includegraphics[width=0.5\linewidth]{logo-large-cropped.png}
    \vspace{30pt}
    \begin{tabular}{c}
      Supervised by: Prof. Zhiyuan Luo \\
      Department of Computer Science \\
      Royal Holloway University of London
    \end{tabular}
  \end{center}
  \newpage
}

\title{Project Plan: Comparison of Machine Learning Algorithms}
\author{Nathan Kurien}
\date{24/09/2023} % Date? 

\begin{document}

\maketitle

\section{Abstract}
% Your abstract goes here.
% Provide an overview of the aims and objectives for the project.

Machine learning (ML) has taken the world by storm in recent times. In the era of unprecedented data generation, ML has emerged as a pivotal technology with profound implications for industries and everyday life. Its ability to extract meaningful patterns from vast datasets has revolutionized decision-making processes and enhanced automation across diverse sectors.\par
Supervised Learning algorithms exist to solve regression and classification tasks, and this project will focus primarily on classification. Within the ML landscape, classification algorithms play a critical role in tasks such as image recognition, spam filtering, medical diagnosis, and more. \par
Classification problems involve using an algorithm, and a given set of inputs used as trained features, to classify any given input with a predefined label. In other words, Classification problems involve searching for a discrete target value, often among a subset of discrete values, whereas Regression problems involve finding a continuous target value. \par
In this project, the aims will be to focus specifically on Classification problems solved by Supervised Learning, evaluate how they function, and compare how several algorithms perform amongst each other. \par
The first algorithm I'll be exploring will be the Nearest Neighbours algorithm, one of the simplest learning algorithms. I'll be building an implementation of this algorithm along with a Decision Tree algorithm implementation. \par




\section{Milestones}
\subsection{First Term}
\subsubsection{Reports}
% Describe the milestones for the first term, including reports.
Report 1: Nearest Neighbour Algorithms - In this report I hope to discuss the theory behind nearest neighbour algorithms, 1-Nearest Neighbour and K-Nearest Neighbour, and further details behind how they function. \par 
Report 2: Decision Tree Algorithms - In this report I hope the discuss the theory behind the Decision Tree algorithm and explore its measures of uniformity. \par
(Extensions on metrics and SVMs) Ensemble Learning? Kernel Methods?

\subsubsection{Programs}
% Describe the milestones for the first term, including programs.
Proof of Concept Programs in Python

\subsection{Second Term}
\subsubsection{Reports}
% Describe the milestones for the second term, including reports.
[Details of second-term report milestones]

\subsubsection{Programs}
% Describe the milestones for the second term, including programs.
[Details of second-term program milestones]

\section{Risks and Mitigations}
% Discuss potential risks associated with your project and how you plan to mitigate them.
[Discussion of project risks and mitigations]

\section{Bibliography}
% List the sources that helped you decide on your project plan.
\begin{itemize}
  \item Source 1
  \item Source 2
  \item Source 3
  % Add more sources as needed
\end{itemize}

\section{Planning and Time-Scales}
% Provide details on the planning and time-scales for your project.
[Project planning and time-scales]

\end{document}
